\scanpage{98}
% archive foliation: 97
\begin{center}- 13 -\end{center}
\hfill n\textsuperscript{o} 285

formée de sous-groupes distingués fermés de $G$, et qui converge vers $e$; en outre, lorsque les conditions de la prop.24 sont satisfaites, $G_\alpha$ est isomorphe à $G/N_\alpha$. Inversement :

PROPOSITION 25 - \uline{Soient} $G$ \uline{un groupe topologique séparé}, $(H_\alpha)$ \uline{une base de filtre sur} $E$, \uline{formée de sous-groupes distingués fermés et convergente vers} $e$. \uline{Pour} $H_\beta\subset H_\alpha$, \uline{soit} $f_{\alpha\beta}$ \uline{l'application canonique de} $G/H_\beta$ \uline{sur} $G/H_\alpha$ ; \uline{alors} $(G/H_\alpha,f_{\alpha\beta})$ \uline{est un système projectif de groupes topologiques vérifiant les conditions de la prop.24}, \uline{et si} $g_\alpha$ \uline{est l'application canonique de} $G$ \uline{sur} $G/H_\alpha$, $g=(g_\alpha)$ \uline{est un isomorphisme de} $G$ \uline{sur un sous-groupe partout dense de} $G'=\varprojlim G/H_\alpha$.

Pour $H_\beta\subset H_\alpha$,(relation de préordre que l'on notera $\alpha\leq\beta$), $G/H_\alpha$ s'identifie canoniquement à $(G/H_\beta)/(H_\alpha/H_\beta)$ (prop.16), donc $f_{\alpha\beta}$ est un homomorphisme, et il est immédiat que si $\alpha\leq\beta\leq\gamma$, $f_{\alpha\beta}\circ f_{\beta\gamma}=f_{\alpha\gamma}$. En outre on a, pour $H_\beta\subset H_\alpha$, $g_\alpha=f_{\alpha\beta}\circ g_\beta$, donc $g$ est une représentation continue de $G$ dans $G'$ (prop.20), et comme l'intersection des $H_\alpha$ est réduite à $e$, $g$ est injective. Comme $g_\alpha$ est surjective, il en est de même de l'application canonique de $G'$ sur $G/H_\alpha$; la prop.24 montre en outre que $g(G)$ est partout dense dans $G'$. Enfin, soit $V$ un voisinage de $e$ dans $G$, et soit $\alpha$ un indice tel que $H_\alpha\subset V$ ; l'ensemble des $y'=g(y)$ tels que $g_\alpha(y)\in g_\alpha(V)$ est un voisinage $W$ de l'élément neutre dans le sous-groupe $g(G)$ de $G'$. Mais la relation $g_\alpha(y)\in g_\alpha(V)$ entraîne par définition $y\in VH_\alpha\subset V^2$ ; on voit donc que $g(V^2)$ est un voisinage de l'élément neutre dans $g(G)$, ce qui achève la démonstration.

On dit que $g$ est \uline{l'application canonique} de $G$ dans $G'=\varprojlim G/H_\alpha$.

On notera que l'on n'a pas nécessairement $g(G)=G'$ (exerc).
